% !TeX root = ../../Book2.tex
\chapter{弱解与 Lax–Milgram 方法}
\label{b2:ch:15}
\BookChapterNumber{b2:ch:15}
% 本章摘要（不计入编号）
\begin{chapterabstract}
本章把散度型方程写成 Hilbert 空间中的测试等式。先明确 Sobolev 边界条件与负阶源项，再完整证明 Lax--Milgram 定理，并以 Poisson 方程说明存在性、唯一性、稳定性和有限维近似。一般有界可测系数无需具有经典导数；一致椭圆性控制主部，Gårding 不等式处理低阶项。最后借紧嵌入与第一卷的 Fredholm 理论，刻画原方程可能出现的有限维核及伴随相容条件。
\end{chapterabstract}

% 本章目标（不计入编号）
\begin{learninggoals}
\begin{itemize}
\item\label{b2:item:15-goal-weak}辨认测试空间所编码的内部方程与边界条件，正确使用 \(H^{-1}\) 数据。
\item\label{b2:item:15-goal-lax}分别核验型的有界性、强制性与空间完备性，掌握 Lax--Milgram 的闭值域证明。
\item\label{b2:item:15-goal-stability}用能量等式证明唯一性、比较和数据稳定性，区分 Galerkin 最佳与准最佳逼近，并处理非齐次边界延拓。
\item\label{b2:item:15-goal-general}说明有界可测系数已足以给出弱存在性，区别原问题与移位问题。
\item\label{b2:item:15-goal-fredholm}在选读范围内，用伴随核写出 Fredholm 相容条件。
\end{itemize}
\end{learninggoals}

设一根杆的导热系数在接头处突然变化。平衡温度不应跳跃，穿过接头的热通量也应保持一致，但温度斜率可以不同。若只寻找处处二次可微的解，便会把这样的基本模型排除在外。弱形式保留可积梯度和通量平衡，把经典导数无法直接表达的问题放进一个完备的函数空间。

\ProseParagraphBreak

上一章用极小化构造对称能量的解；本章直接求解双线性等式，从而允许非对称系数。Hilbert 空间表示与紧算子的抽象结论采用第一卷第八章和附录 C 的证明，Sobolev 边界与紧嵌入采用第一卷第二十二至二十五章。有关多维 Sobolev 空间与椭圆边值问题的变分表述，可参阅 Haim Brezis 的《Functional Analysis, Sobolev Spaces and Partial Differential Equations》\cite[Chapter 9]{Brezis2011Functional}；Lax--Milgram、能量估计与 Fredholm 择一的进一步阅读见 Lawrence C. Evans 的《Partial Differential Equations》\cite[Section 6.2]{Evans2010Partial}。本章将先完成这些结果在所用空间中的证明，再讨论有限维近似与粗糙系数；进一步阅读不承担正文缺失的证明步骤。

% !TeX root = ../../Book2.tex
\section{弱形式}
\label{b2:sec:1501}

分段材料中的温度可以连续而斜率发生跳跃；此时逐点写二阶导数会掩盖界面上的通量条件。弱形式把微分移到测试函数上，使未知函数只需具有有限能量，同时仍保留方程及界面平衡。本节将第一章的局部弱解概念与上一章的 Sobolev 边界类结合起来，并为比 \(L^2\) 更粗糙的源项确定合适空间。

\subsection{测试函数}
\label{b2:subsec:150101}

在开集 \(\Omega\subset\mathbb R^d\) 中，测试函数取自 \(C_c^\infty(\Omega)\)。它们的支集离边界有正距离，因此分部积分不会引入尚未定义的边界值。若 \(u\in H_{\mathrm{loc}}^1(\Omega)\)、\(A\in L_{\mathrm{loc}}^\infty(\Omega;\mathbb R^{d\times d})\)，则
\[\langle-\operatorname{div}(A\nabla u),\varphi\rangle
 =\int_\Omega A\nabla u\cdot\nabla\varphi\,\mathrm dx\]
定义了一个分布。矩阵约定为 \((A\nabla u)_i=\sum_j a_{ij}\partial_j u\)，因而散度作用于通量的第一个指标。

\ProseParagraphBreak

局部测试函数足以辨认内部方程，但能量估计常须取 \(\varphi=u\)、截断后的 \(u\) 或差商。它们通常不光滑，也不紧支集。合法性依靠的是密度与连续延拓：若左右两边对 \(H_0^1\) 范数连续，而 \(C_c^\infty\) 在该空间稠密，则测试等式自动推广到整个 \(H_0^1\)。若只知道局部 \(H^1\)，便先乘紧支集截断函数，在较小区域完成同样步骤。不能在未建立这些连续性之前直接“取测试函数为解”。

\subsection{弱解}
\label{b2:subsec:150102}

以下主要讨论实函数；实系数方程的复值版本可分别作用于实部与虚部。设整数 \(d\geq1\)，\(\Omega\subset\mathbb R^d\) 为开集，\(A\in L^\infty(\Omega;\mathbb R^{d\times d})\)，\(f\in\mathcal D'(\Omega;\mathbb R)\)。本章主要研究有限全局能量的弱解，因此称 \(u\in H^1(\Omega;\mathbb R)\) 满足散度方程的弱形式，是指
\[\int_\Omega A\nabla u\cdot\nabla\varphi=\langle f,\varphi\rangle
 \qquad(\varphi\in C_c^\infty(\Omega)).\]
第一章的局部弱解只要求 \(u\in H_{\mathrm{loc}}^1(\Omega)\)，使用的是同一个测试恒等式；本处的 \(H^1(\Omega)\) 要求更强，还控制整个区域上的函数与梯度平方积分。全局能量弱解必为局部弱解，反向一般不成立。两者都尚未包含边界条件；零 Dirichlet 条件还要在下一小节另行加入。对于 \(A\in C^1(\Omega;\mathbb R^{d\times d})\)、\(u\in C^2(\Omega)\) 及连续函数 \(f\)，分部积分和测试函数基本引理使该恒等式与经典方程等价；粗糙情形则以积分等式本身为准。

\ProseParagraphBreak

这个空间差别已经出现在三维的去心球上。取
\[
 \Omega=B(0,1)\setminus\{0\}\subset\mathbb R^3,\qquad u(x)=|x|^{-1}.
\]
函数在 \(\Omega\) 内光滑；径向公式给出
\(\Delta u=u''(r)+2u'(r)/r=2r^{-3}-2r^{-3}=0\)。
每个紧含于 \(\Omega\) 的子集都离原点有正距离，所以 \(u\in H_{\mathrm{loc}}^1(\Omega)\)，并对所有紧支集测试函数满足 \(-\Delta u=0\)。但球坐标计算给出
\[
 \begin{aligned}
 \int_\Omega |u|^2\,\mathrm dx&=4\pi\int_0^1\mathrm dr=4\pi,\\
 \int_{\{\varepsilon<|x|<1\}}|\nabla u|^2\,\mathrm dx
 &=4\pi\int_\varepsilon^1r^{-2}\,\mathrm dr\\
 &=4\pi(\varepsilon^{-1}-1)\longrightarrow+\infty.
 \end{aligned}
\]
故它是局部弱解，却不属于本章的全局能量解类。这里原点已从区域删去，测试函数接触不到原点；不能把在整个球内的点源方程与去心球上的齐次方程混同。

\Needspace{6\baselineskip}
\begin{proposition}[一维界面的通量条件]
\label{b2:prop:15-interface-flux}
设 \(-1<s<1\)，实系数 \(a\) 在 \((-1,s)\)、\((s,1)\) 上分别为正常数，实函数 \(u\) 在两闭子区间上分别为 \(C^2\)，且在 \(s\) 处连续。若 \(-(au')'=f\) 在两侧经典成立，其中实值 \(f\in L^2(-1,1)\)，则 \(u\in H^1(-1,1)\) 在整个区间满足 \(-(au')'=f\) 的弱形式，当且仅当
\[a(s-)u'(s-)=a(s+)u'(s+).\]
\end{proposition}
\begin{proof}
分段光滑且连续的 \(u\) 属于 \(H^1\)，其弱导数就是分段的一阶导数。在两段分别积分，外端因测试函数紧支集而消失，内端得到
\[\int_{-1}^1 au'\varphi'=\int_{-1}^1 f\varphi
  +\bigl(a(s-)u'(s-)-a(s+)u'(s+)\bigr)\varphi(s).\]
测试函数在 \(s\) 处可取任意值，因此等式对全部测试函数成立恰当且仅当括号为零。
\end{proof}

例如 \(a=1\) 于左半区间、\(a=2\) 于右半区间，令 \(u(x)=x\)（\(x\le0\)）、\(u(x)=x/2\)（\(x\ge0\)）。其导数跳跃，但通量恒为 \(1\)，所以 \(-(au')'=0\) 弱成立。若误要求 \(u'\) 连续，就会排除这个正确的平衡解；若允许 \(u\) 本身跳跃，又会产生分布导数中的点质量，通常已不属于所选能量空间。

\subsection{Sobolev 边界条件}
\label{b2:subsec:150103}

对任意有界开集，齐次 Dirichlet 条件定义为 \(u\in H_0^1(\Omega)\)；非齐次数据由一个延拓 \(g\in H^1(\Omega)\) 表示，条件为 \(u-g\in H_0^1(\Omega)\)。第一卷定理 23.1.6 说明，在有界 Lipschitz 区域上，这恰好等价于 Sobolev 迹满足 \(\operatorname{Tr}u=\operatorname{Tr}g\)。粗糙区域上仍可使用闭包定义，但不能未经证明把它等同于某个朴素逐点边界条件。

\ProseParagraphBreak

自然边界条件的写法不同。以非空有界 \(C^1\) 开集 \(\Omega\subset\mathbb R^d\) 为例，若实函数 \(u\) 在 \(\overline\Omega\) 的某开邻域内为 \(C^2\)，记 \(f=-\Delta u\)、\(g_N=\partial_\nu u\)，则在 \(H^1(\Omega;\mathbb R)\) 的全部方向测试时，Green 公式是
\[\int_\Omega\nabla u\cdot\nabla v
  =\int_\Omega fv+\int_{\partial\Omega}g_N\operatorname{Tr}v\,\mathrm dS.\]
这条公式先对光滑 \(v\) 成立，再由 \(H^1\) 光滑逼近和迹算子的连续性推广到全部 \(H^1\) 测试函数。它编码 \(\partial_\nu u=g_N\)，而不是要求 \(u\) 在边界为零。取 \(v=1\) 得到必要相容条件 \(\int_\Omega f+\int_{\partial\Omega}g_N=0\)；若区域不连通，还须在每个连通分支上分别成立，见\cref{b2:prop:13-neumann-compatibility}。测试空间不同，编码的边界问题便不同；仅看内部方程无法辨认二者。

\ProseParagraphBreak

即使区域光滑，也不宜把任意边界函数都当作合法 \(H^1\) Dirichlet 数据。数据应当处于迹空间，并存在 \(H^1\) 延拓；第一卷第二十三章在 Lipschitz 图中用分数阶范数刻画这种要求。给定一个 \(H^1\) 延拓 \(g\)，就可以把满足该边界数据的函数写为 \(g+H_0^1(\Omega)\) 中的元素。

\subsection{负阶 Sobolev 空间初步}
\label{b2:subsec:150104}

设整数 \(d\geq1\)，\(\Omega\subset\mathbb R^d\) 是非空有界开集，以梯度内积
\(\langle u,v\rangle_{H_0^1}=\int_\Omega\nabla u\cdot\nabla v\)
赋予实空间 \(H_0^1(\Omega;\mathbb R)\) Hilbert 结构。第一卷定理24.1.2的零边界 Poincaré 不等式给出
\(\|v\|_2\leq C_P\|\nabla v\|_2\)，所以梯度范数与完整 \(H^1\) 范数等价，这个空间仍完备。定义
\[H^{-1}(\Omega)=(H_0^1(\Omega))^*,\qquad
 \lVert f\rVert_{H^{-1}}=\sup_{\lVert\nabla v\rVert_2\le1}|\langle f,v\rangle|.
\]
这是本章使用的 \term[fujie-sobolev-kongjian]{负阶 Sobolev 空间}[negative-order Sobolev space]的第一种情形。换成全 \(H^1\) 范数会得到等价的对偶范数；连续对偶中的元素不变，但下面的精确范数等式使用的是梯度范数约定。每个对偶元素 \(f\in H^{-1}(\Omega)\) 限制到 \(C_c^\infty(\Omega)\) 都给出分布：对支集落在固定紧集 \(K\Subset\Omega\) 内的 \(\varphi\)，有
\[
 |\langle f,\varphi\rangle|
 \leq\|f\|_{H^{-1}}\,|K|^{1/2}\sup_K|\nabla\varphi|,
\]
所以配对在测试函数拓扑下连续。密度保证这个限制是单射，因此也可以把 \(H^{-1}\) 看作一类分布。

实值函数 \(f\in L^2(\Omega)\) 通过 \(\langle f,v\rangle=\int_\Omega fv\) 进入这个对偶，因为
\[
 |\langle f,v\rangle|\leq \|f\|_2\|v\|_2
                 \leq C_P\|f\|_2\|\nabla v\|_2.
\]
更一般地，若 \(F\in L^2(\Omega;\mathbb R^d)\)，则
\(\langle-\operatorname{div}F,v\rangle=\int_\Omega F\cdot\nabla v\)
在 \(H_0^1\) 上连续，范数不超过 \(\|F\|_2\)。后一个例子并不要求散度本身是一项 \(L^2\) 函数；它只要求给出的测试作用有能量范数上界。

\Needspace{6\baselineskip}
\begin{proposition}[负阶数据的通量表示]
\label{b2:prop:15-negative-representation}
设整数 \(d\geq1\)，\(\Omega\subset\mathbb R^d\) 为非空有界开集，\(H_0^1(\Omega;\mathbb R)\) 取梯度内积，\(H^{-1}(\Omega)=(H_0^1(\Omega;\mathbb R))^*\) 取对偶范数
\(\|f\|_{H^{-1}}=\sup_{\|\nabla v\|_2\leq1}|\langle f,v\rangle|\)。
对每个 \(f\in H^{-1}(\Omega)\)，存在唯一 \(w\in H_0^1(\Omega;\mathbb R)\) 使
\[
 \langle f,v\rangle=\int_\Omega\nabla w\cdot\nabla v\,\mathrm dx
 \qquad(v\in H_0^1(\Omega;\mathbb R)).
\]
因而 \(F_0=\nabla w\in L^2(\Omega;\mathbb R^d)\) 是 \(f=-\operatorname{div}F_0\) 的通量表示，且
\(\|F_0\|_2=\|f\|_{H^{-1}}\)。
任一其他通量表示恰为 \(F=F_0+G\)，其中 \(G\in L^2(\Omega;\mathbb R^d)\)、\(\operatorname{div}G=0\) 于 \(\mathcal D'(\Omega)\)，并满足
\[
 \|F\|_2^2=\|F_0\|_2^2+\|G\|_2^2.
\]
特别地，\(F_0\) 是所有通量表示中唯一的最小 \(L^2\) 范数表示。
\end{proposition}
\begin{proof}
Poincaré 不等式使梯度内积给出完备的 Hilbert 空间。直接在这个空间上应用第一卷定理8.4.7的 Riesz 表示，得到唯一 \(w\) 及所列配对，且
\[
 \|f\|_{H^{-1}}=\|w\|_{H_0^1}=\|\nabla w\|_2.
\]
取 \(F_0=\nabla w\) 就得到通量表示。

若 \(F\) 也表示 \(f\)，则 \(G=F-F_0\) 满足
\(\int_\Omega G\cdot\nabla v=0\) 对全部 \(v\in H_0^1\) 成立。特别取紧支集光滑 \(v\)，便知 \(\operatorname{div}G=0\) 于分布意义。反过来，若 \(G\in L^2\) 的分布散度为零，等式先对紧支集测试函数成立；由于其左端对梯度范数连续，密度使它对全部 \(H_0^1\) 成立，所以 \(F_0+G\) 仍表示 \(f\)。最后取 \(v=w\)，有 \(\int_\Omega G\cdot F_0=0\)，展开平方即得范数恒等式；等号 \(\|F\|_2=\|F_0\|_2\) 当且仅当 \(G=0\)，证明最小范数表示的唯一性。
\end{proof}

因此负阶数据确定的是通量的散度；最小范数准则才从所有通量中选出梯度通量。可以自由叠加的 \(G\) 与全部零边界梯度正交，但这个条件只等价于分布散度为零，不能在未定义法向迹时再解释为逐点零法向通量。

这个空间也确实容纳比函数更粗糙的数据，但是否容纳点源取决于维数。以下直接检查点值泛函；所需的不是“点质量是分布”，而是它对能量范数连续。

\Needspace{6\baselineskip}
\begin{proposition}[点质量与能量对偶]\label{b2:prop:15-point-mass-dimension}
设整数 \(d\geq1\)，\(\Omega\subset\mathbb R^d\) 为非空有界开集，\(H^{-1}(\Omega)\) 表示实空间 \(H_0^1(\Omega;\mathbb R)\) 在梯度范数下的连续对偶。对 \(x_0\in\Omega\)，点质量分布定义为 \(\langle\delta_{x_0},\varphi\rangle=\varphi(x_0)\)、\(\varphi\in C_c^\infty(\Omega)\)。若 \(d=1\)、\(\Omega=I=(a,b)\) 为有限开区间，则 \(\delta_{x_0}\in H^{-1}(I)\)；若 \(d\geq2\)，则 \(\delta_{x_0}\notin H^{-1}(\Omega)\)。
\end{proposition}
\begin{proof}
在 \(I=(a,b)\) 上，对光滑零边界函数由
\(\varphi(x_0)=\int_a^{x_0}\varphi'\) 得
\(|\varphi(x_0)|\leq\sqrt{b-a}\|\varphi'\|_2\)。密度把点值泛函延拓到 \(H_0^1(I)\)，其值与绝对连续代表的点值一致：第一卷命题21.4.1给出一维 \(H^1\) 函数的绝对连续代表，而上述增量估计保证零边界光滑逼近的点值趋于该代表的点值。

高维先平移并缩放一个紧含于 \(\Omega\) 的球，使观察点为零、球半径为一；这个固定换元只改变范数常数。若 \(d>2\)，取 \(\eta\in C_c^\infty(B(0,1))\)、\(\eta(0)=1\)，令 \(\eta_\varepsilon(x)=\eta(x/\varepsilon)\)。其点值为一，梯度范数平方却是 \(\varepsilon^{d-2}\|\nabla\eta\|_2^2\to0\)。不存在同时控制这些点值的能量常数。

二维的简单缩放使梯度范数不变，须改用较缓的对数过渡。取 \(0<\varepsilon<1/2\)，令径向函数 \(w_\varepsilon\) 在 \(r\leq\varepsilon^2\) 等于一，在 \(\varepsilon^2<r<\varepsilon\) 等于 \(\log(\varepsilon/r)/\log(1/\varepsilon)\)，在 \(r\geq\varepsilon\) 等于零。它连续、分段光滑，且
\[
 \|\nabla w_\varepsilon\|_2^2
 =\frac{2\pi}{\log(1/\varepsilon)^2}
   \int_{\varepsilon^2}^{\varepsilon}\frac{\mathrm dr}{r}
 =\frac{2\pi}{\log(1/\varepsilon)}\longrightarrow0.
\]
以非负单位质量磨光核在半径小于 \(\varepsilon^2/2\) 的尺度卷积，得到光滑紧支集函数：中心邻域仍恒为一，支集仍在单位球内，Young 不等式使梯度范数不增。因此光滑测试函数同样违反 \(|\varphi(0)|\leq C\|\nabla\varphi\|_2\)。这正是否定点质量属于对偶所需的结论。
\end{proof}
一维点源的精确范数将在习题中进一步计算；高维的替代训练则计算小球平均源的能量如何发散。\Cref{b2:tab:15-source-classes}把本节已证的源项范围放在一起，便于应用后面的存在性定理前先核对数据。

\begin{table}[htbp]
\centering
\small
\caption[四类源项与能量对偶]{四类源项与能量对偶。固定整数 \(d\geq1\) 和非空有界开集 \(\Omega\subset\mathbb R^d\)，所有空间取实数域，\(H^{-1}(\Omega)=(H_0^1(\Omega;\mathbb R))^*\) 使用梯度对偶范数，\(C_P\) 为 \(\Omega\) 的零边界 Poincaré 常数。点质量两行另按所列维数和位置取数据；结论均已在正文证明。}
\label{b2:tab:15-source-classes}
\begin{tabularx}{\linewidth}{@{}>{\raggedright\arraybackslash}p{3.8cm}>{\centering\arraybackslash}p{1.5cm}>{\raggedright\arraybackslash}X@{}}
\toprule
数据形式及条件 & 属于 \(H^{-1}\) & 连续性估计或障碍 \\
\midrule
\(f\in L^2(\Omega)\) & 是 & \(\|f\|_{H^{-1}}\leq C_P\|f\|_2\)，使用零边界 Poincaré 不等式 \\
\(-\operatorname{div}F\)，\(F\in L^2(\Omega;\mathbb R^d)\) & 是 & \(\|-\operatorname{div}F\|_{H^{-1}}\leq\|F\|_2\)；不要求散度是普通函数 \\
\(d=1\)，\(\Omega=(a,b)\)，\(a<s<b\) 时的 \(\delta_s\) & 是 & \(|v(s)|\leq\sqrt{b-a}\|v'\|_2\)，点值取绝对连续代表 \\
\(d\geq2\)，\(x_0\in\Omega\) 时的 \(\delta_{x_0}\) & 否 & 测试函数点值固定为一，梯度能量可趋零；高于二维用缩放，二维用对数截断 \\
\bottomrule
\end{tabularx}
\end{table}

后面的零 Dirichlet 弱存在性定理只适用于已核验属于 \(H^{-1}\) 的数据，不能把任意点源自动代入；点质量的分布意义本身并不提供能量对偶中的连续性。

% !TeX root = ../../Book2.tex
\section{弱 \texorpdfstring{\(L^p\)}{Lᵖ} 空间}
\label{b2:sec:1702}

在端点问题中，函数超过高度 \(t\) 的区域往往比函数范数更容易估计。本节把高度与区域大小之间的关系写成函数空间语言。底空间可以是一般测度空间，重排部分以无原子的 Lebesgue 空间为例。

\subsection{分布函数}
\label{b2:subsec:170201}

设 \(f\) 可测，其\term[fenbu-hanshu]{分布函数}[distribution function]定义为 \(d_f(t)\coloneqq\mu\{|f|>t\}\)，\(t>0\)。它单调不增；对递减高度 \(t_n\downarrow t\)，集合 \(\{|f|>t_n\}\) 递增到 \(\{|f|>t\}\)，所以 \(d_f(t_n)\to d_f(t)\)。严格不等号保证这里的右连续性。
\symidx{\(d_f(t)\)：函数的分布函数}
\begin{proposition}[层蛋糕公式]\label{b2:thm:17-layer-cake}
对 \(0<p<\infty\) 和可测 \(f\)，有允许无穷值的恒等式
\[
 \int|f|^p\,\mathrm d\mu=p\int_0^\infty t^{p-1}d_f(t)\,\mathrm dt.
\]
\end{proposition}
\begin{proof}
逐点使用 \(|f(x)|^p=\int_0^\infty p t^{p-1}\mathbf1_{\{t<|f(x)|\}}\,\mathrm dt\)，再用 Tonelli 换序。非负性使此式在两端都无穷时仍有明确意义。
\end{proof}
在有限测度空间上，常函数的分布函数是一个阶梯函数；在 \(\mathbb R^d\) 上，\(|x|^{-d/p}\mathbf1_{\{|x|<1\}}\) 的高水平尾部则正好具有 \(t^{-p}\) 的尺度。这两类对象分别说明有界函数与临界奇性的分布图形。

\subsection{弱型 \texorpdfstring{\(L^p\)}{Lp}}
\label{b2:subsec:170202}

对 \(0<p<\infty\)，若
\[
 \|f\|_{p,\infty}\coloneqq\sup_{t>0}t\cdot d_f(t)^{1/p}<\infty,
\]
则称 \(f\) 属于\term[ruo-lp-kongjian]{弱 \(L^p\) 空间}[weak Lp space]，记为 \(L^{p,\infty}\)。这里给出的是拟范数：由 \(\{|f+g|>t\}\subset\{|f|>t/2\}\cup\{|g|>t/2\}\) 得到带常数的三角不等式，不宣称常数为一。
\begin{proposition}\label{b2:thm:17-weak-inclusions}
有 \(\|f\|_{p,\infty}\leq\|f\|_p\)。若 \(p>1\)、\(\mu(E)<\infty\)，则
\[
 \int_E|f|\,\mathrm d\mu\leq\dfrac{p}{p-1}\|f\|_{p,\infty}\cdot\mu(E)^{1-1/p}.
\]
反之，在半有限测度空间中，若所有有限测度 \(E\) 都满足右式中某个固定常数的估计，则 \(f\in L^{p,\infty}\)。
\end{proposition}
\begin{proof}
第一项由 Chebyshev 得到。记 \(A=\|f\|_{p,\infty}\)、\(m=\mu(E)\)。在 \(m=0\) 或 \(A=0\) 时结论直接成立；否则按层蛋糕公式，积分不超过 \(\int_0^\infty\min\{m,A^pt^{-p}\}\,\mathrm dt\)。在 \(t=A m^{-1/p}\) 分开积分即得。反向在 \(E=\{|f|>t\}\) 上使用 \(t\mu(E)\leq\int_E|f|\)，若集合测度无穷，先取其中有限测度子集；在半有限空间上便推出 \(t\mu(E)^{1/p}\leq C\)。本卷的 Lebesgue 空间满足这一条件。
\end{proof}
积分集合刻画解释了为什么 \(p>1\) 时弱空间中的函数仍局部可积，而 \(p=1\) 时不能作同样断言。

\ProseParagraphBreak

若算子 \(T\) 满足 \(\|Tf\|_{q,\infty}\leq A\|f\|_p\)，称它为\term[ruoxing-guji]{弱型估计}[weak type estimate]的 \((p,q)\) 版本；强型把目标换为 \(L^q\)。当 \(q=\infty\) 时这里直接使用本性上确界，不另定义分布函数拟范数。

% 扩写定位：b2:subsec:170202
在 Lebesgue 空间上，集合积分刻画还给出一个真正的范数。设 \(1<p<\infty\)，定义
\[
 N_p(f)\coloneqq\sup_{0<|E|<\infty}|E|^{-1/p'}\int_E|f|.
\]
这里可以只取有界可测集合：一般集合与球相交后逐步逼近即可。由上一命题，
\[
 \|f\|_{p,\infty}\leq N_p(f)\leq p'\|f\|_{p,\infty}.
\]
左侧来自超水平集合及其有限测度子集。绝对齐次性、三角不等式和 \(N_p(f)=0\) 蕴含 \(f=0\) 几乎处都直接由定义得到。因而在这个指数范围内，弱空间可以赋予与原拟范数等价的 Banach 范数。
\begin{proposition}\label{b2:prop:17-weak-complete}
\(L^{p,\infty}(\mathbb R^d)\) 关于 \(N_p\) 完备。若 \(k\in L^1\)，则卷积具有连续延拓，并满足
\[
 N_p(k*f)\leq\|k\|_1N_p(f).
\]
\end{proposition}
\begin{proof}
设 \(f_n\) 在 \(N_p\) 中 Cauchy。在每个有限球 \(B\) 上，
\(\|f_n-f_m\|_{L^1(B)}\leq |B|^{1/p'}N_p(f_n-f_m)\)，故得到相容的局部 \(L^1\) 极限 \(f\)。取沿所有整数球几乎处收敛的对角子列。对任意有限测度集合 \(E\)，Fatou 给出
\(\int_E|f_n-f|\leq\liminf_m\int_E|f_n-f_m|\)；先固定足够大的 \(n\)，再对 \(E\) 取上确界，得到 \(N_p(f_n-f)\to0\)。这同时证明 \(f\) 的范数有限。

对卷积，Tonelli 和平移不变性给出
\[
 \int_E\int |k(y)|\,|f(x-y)|\,\mathrm dy\,\mathrm dx
 \leq |E|^{1/p'}\|k\|_1N_p(f).
\]
右端有限，故积分在每个有限球上几乎处绝对收敛；其绝对值控制 \(|k*f|\)，再取上确界便得结论。
\end{proof}
例如，临界势核的弱范数可以在与一个可积平滑核卷积后保持同一指数。这里并不需要卷积核正定，也不要求 \(f\) 本身可积；用局部集合积分而不是全空间 \(L^1\) 积分，正好避免了这一障碍。

\ProseParagraphBreak

完备性不能误读成平滑紧支集函数在弱范数下稠密。令 \(f(x)=|x|^{-d/p}\mathbf1_{\{|x|<1\}}\)，并令 \(g\) 有界。在足够小的球上，\(|f-g|\geq |f|/2\)；高水平集合的尺度于是给出 \(\|f-g\|_{p,\infty}\geq c_{d,p}>0\)。有界截断 \(f\mathbf1_{\{|f|\leq N\}}\) 因此也不收敛到 \(f\) 的弱范数。处理算子的弱空间延拓时，需要使用已证明的积分定义，或者指定额外的收敛方式，不能照搬强 \(L^p\) 的稠密性论证。

\subsection{Lorentz 空间 初步}
\label{b2:subsec:170203}

定义递减重排 \(f^*(s)\coloneqq\inf\{t\geq0:d_f(t)\leq s\}\)，\(s>0\)，下确界为空时取无穷。它把函数值按大小重新排列，而不保留位置。\term[lorentz-kongjian]{Lorentz 空间}[Lorentz space]通过
\[
 \|f\|_{L^{p,q}}\coloneqq
 \left(\int_0^\infty \bigl(s^{1/p}f^*(s)\bigr)^q\dfrac{\mathrm ds}{s}\right)^{1/q}
 \quad(0<q<\infty)
\]
衡量整个高度分布，\(q=\infty\) 时改为上确界。这里仍用拟范数约定。重排与原函数具有相同分布，故层蛋糕公式给出 \(\|f\|_{L^{p,p}}=\|f\|_p\)，而 \(L^{p,\infty}\) 与上段定义一致。
\symidx{\(f^*(s)\)：递减重排}

\ProseParagraphBreak

对 \(1<p<\infty\)，若 \(f\in L^{p,1}\)、\(g\in L^{p',\infty}\)，则
\[
 \int|fg|\leq\|f\|_{L^{p,1}}\cdot\|g\|_{L^{p',\infty}}.
\]
确实，写 \(f=\int_0^\infty\mathbf1_{\{|f|>t\}}\,\mathrm dt\)，集合积分的重排界为 \(\int_E|g|\leq\int_0^{|E|}g^*(s)\,\mathrm ds\)。这个界由层蛋糕和 \( |E\cap\{|g|>u\}|\leq\min\{|E|,d_g(u)\}\) 得到。换序后得到 \(\int|fg|\leq\int_0^\infty f^*(s)g^*(s)\,\mathrm ds\)，再代入 \(g^*(s)\leq\|g\|_{p',\infty}s^{-1/p'}\)。这是把一个临界弱估计与更细的输入空间配对的基本方式。

\ProseParagraphBreak

例如令 \(f^*(s)\) 在充分小的 \(s>0\) 上等于 \(s^{-1/p}\bigl(\log(e/s)\bigr)^{-\beta}\)，其余部分取单调有限支集延拓。其 \(L^{p,q}\) 积分在 \(s=0\) 收敛当且仅当 \(\beta q>1\)，弱 \(L^p\) 只要求 \(\beta\geq0\)。第二个指数记录了强空间 \(L^p\) 所不能区分的端点细节。

% 扩写定位：b2:subsec:170203
几种分布尾部可以帮助区分这些空间。首先，函数在有限测度集合上有界，并不能说明其重排在大 \(s\) 上怎样；当函数支集也有限时，重排在支集测度之后为零。其次，一个幂次奇性在小球上产生高水平的幂次尾部，决定局部的可积临界；同一幂次放在球外，则产生低水平的大测度尾部，决定远场的可积临界。两种限制可以落在相反方向，因此全空间上的空间包含关系不能只看一个局部模型。

\ProseParagraphBreak

具体地，若 \(f(x)=|x|^{-a}\mathbf1_{\{|x|<1\}}\)，径向积分说明 \(f\in L^p\) 当且仅当 \(ap<d\)，临界 \(ap=d\) 只有弱 \(L^p\)。若 \(g(x)=|x|^{-a}\mathbf1_{\{|x|>1\}}\)，则 \(g\in L^p\) 当且仅当 \(ap>d\)，临界同样只有弱空间。把两部分加在一起得到纯幂次 \(|x|^{-a}\)；它属于临界弱空间，却不属于任何全空间强 \(L^p\)。这不是弱空间定义的反常现象，而是它准确记录了在每个高度上超水平球的尺度。

\ProseParagraphBreak

有时输入的第二指数需要比弱空间更严格。设 \(f^*(s)\) 在 \(s\downarrow0\) 时具有下列量级：
\[
 f^*(s)\asymp s^{-1/p}(\log(e/s))^{-\beta}.
\]
在其余部分作单调的有限支集延拓。代入 Lorentz 定义，换元 \(v=\log(e/s)\)，所得尾积分为 \(\int^\infty v^{-\beta q}\,\mathrm dv\)。因此固定 \(p\) 后，\(q\) 越小要求越强；例如在 \(\beta=1/p\) 时恰好不属于 \(L^p\)，但可能属于具有更大第二指数的空间。这里说的是原点附近的量级，若 \(\beta\) 大，原显示公式在整个区间上未必递减，应以这种局部模型加单调延拓来实现重排。

\ProseParagraphBreak

前述端点配对中，\(L^{p,1}\) 一般不能放宽为 \(L^p\)。取临界重排
\[
 g^*(s)=s^{-1/p'}\mathbf1_{\{0<s<1\}}.
\]
它属于弱 \(L^{p'}\)。若 \(f^*(s)=s^{-1/p}(\log(e/s))^{-\beta}\) 且 \(1/p<\beta\leq1\)，则 \(f\in L^p\)，但同位置递减模型的乘积积分归结为
\[
 \int_0^1s^{-1}(\log(e/s))^{-\beta}\,\mathrm ds=\infty.
\]
取 \(\beta>1\) 则 \(f\in L^{p,1}\)，前面证明的配对界保证收敛。于是把 \(L^{p,1}\) 换成 \(L^p\) 一般会破坏结论，这个反例直接指出第二指数在临界积分中的作用。

\ProseParagraphBreak

在有限测度区域 \(E\) 上，若 \(0<q<p\)，同样按层蛋糕把高度在 \(\|f\|_{p,\infty}|E|^{-1/p}\) 处分开，可得
\[
 \|f\|_{L^q(E)}
 \leq\left(\frac p{p-q}\right)^{1/q}|E|^{1/q-1/p}\|f\|_{p,\infty}.
\]
这不是弱 \(L^p\) 的强端点；它降低了目标指数，同时用到了区域测度。若 \(q=p\)，高度积分再次成为对数发散；若 \(|E|=\infty\)，所显示的区域因子不提供有限上界。因而使用这一局部提升时，指数严格性与区域有限性都必须保留。

\ProseParagraphBreak

还有一种较弱的输入端点，只先在集合指示函数上检验算子。若线性算子 \(T\) 在有限支集简单函数上定义，并满足
\(\|T\mathbf1_E\|_{q,\infty}\leq A|E|^{1/p}\)、\(1\leq p<\infty\)、\(1<q<\infty\)，称这是一项限制弱型估计。它并不直接给出全部 \(L^p\) 输入的弱型界，但足以推出
\[
 N_q(Tf)\leq C_{p,q}A\|f\|_{L^{p,1}}.
\]
对非负简单函数，层蛋糕是有限个集合指示函数的正线性组合。范数三角不等式与弱范数等价给出
\[
 N_q(Tf)\leq q'A\int_0^\infty d_f(t)^{1/p}\,\mathrm dt
 =\frac{q'A}{p}\|f\|_{L^{p,1}}.
\]
最后的等式由 Tonelli 对重排的层蛋糕积分得到。一般复简单函数分解为实部与虚部的正负部分，常数只增加一个固定因子。

\ProseParagraphBreak

有限支集简单函数在 \(L^{p,1}\) 中稠密：先截去位置尾部、大值尾部和极小值尾部，重排积分的可积性保证误差趋零，再在剩余有界有限测度区域作简单逼近。因 \(N_q\) 完备，以上估计产生唯一连续延拓。这个结果说明限制弱型端点自然对应更细的 \(L^{p,1}\) 输入，而不能在没有新论证的情况下换成所有 \(L^p\)。当目标 \(q=1\) 时，本段使用的等价 Banach 范数已不存在，仍需另一种处理。

% !TeX root = ../../Book2.tex
\section{Poisson 方程}
\label{b2:sec:1503}

对 Poisson 方程，把二阶导数经分部积分转到测试函数上，就得到关于梯度的双线性等式。Lax--Milgram 定理可以在有界区域的能量空间中求解这一等式，并允许源项属于 \(H^{-1}\)。这种写法将存在性、稳定性与边界数据放在同一个框架中。

本节使用\subsecref{b2:subsec:150104}的实空间约定：\(H^{-1}(\Omega)\) 是带梯度范数的 \(H_0^1(\Omega;\mathbb R)\) 的连续对偶，\(\langle f,v\rangle\) 表示对偶配对。下面常数为一的估计与这一梯度范数约定相配；换用完整 \(H^1\) 范数时，对偶空间不变，数值常数则须作相应比较。

\subsection{弱解存在性}
\label{b2:subsec:150301}

\Needspace{6\baselineskip}
\begin{theorem}[Dirichlet 问题的弱可解性]
\label{b2:thm:15-poisson-existence}
设整数 \(d\geq1\)，\(\Omega\subset\mathbb R^d\) 为非空有界开集，空间取实数域。以梯度范数赋予 \(H_0^1(\Omega)\) 范数，令 \(H^{-1}(\Omega)=(H_0^1(\Omega))^*\)，其对偶范数为 \(\|f\|_{H^{-1}}=\sup_{v\in H_0^1,\,\|\nabla v\|_2\leq1}|\langle f,v\rangle|\)。对每个 \(f\in H^{-1}(\Omega)\)，存在唯一 \(u\in H_0^1(\Omega)\)，使
\[\int_\Omega\nabla u\cdot\nabla v\,\mathrm dx=\langle f,v\rangle
 \qquad(v\in H_0^1(\Omega)).\]
并且 \(\lVert\nabla u\rVert_2\le\lVert f\rVert_{H^{-1}}\)。
\end{theorem}
\begin{proof}
取 \(H=H_0^1(\Omega)\)，内积为 \(\langle u,v\rangle_H=\int\nabla u\cdot\nabla v\)。第一卷定理24.1.2的零边界 Poincaré 不等式适用于任意包含于有限宽条带的开集，故适用于这里的有界 \(\Omega\)，不要求边界正则或连通。它使梯度范数与完整 Sobolev 范数等价，故 \(H\) 完备。型 \(a(u,v)=\langle u,v\rangle_H\) 的有界常数与强制常数均为 \(1\)，右端按 \(H^{-1}\) 定义连续。\cref{b2:thm:15-lax-milgram} 因而给出存在、唯一性及估计。
\end{proof}

本例中 Lax--Milgram 退化为 Riesz 表示；一般系数时，同一个定理才显出更大的作用。该弱解也恰为 \(\tfrac12\lVert\nabla v\rVert_2^2-\langle f,v\rangle\) 的唯一极小元：只要将上一章的二次展开中的 \(L^2\) 配对换成对偶配对，证明逐字成立。两种构造在对称情形相遇，但弱型方法还可以继续处理非对称问题。

\ProseParagraphBreak

对 \(\Omega=(0,1)\)、\(f=\delta_s\)、\(0<s<1\)，解为
\[
 u(x)=\begin{cases}(1-s)x,&0\le x\le s,\\s(1-x),&s\le x\le1.\end{cases}
\]
它连续、分段仿射，导数在 \(s\) 处向下跳 \(1\)。两段分部积分给出 \(\int_0^1u'v'=v(s)\)，先对光滑 \(v\) 成立，再由点值的连续性推广。因此点源产生一个合法弱解，却不产生经典 \(C^2\) 解；这不是存在性定理的缺陷，而是数据本身的粗糙性。

\subsection{唯一性}
\label{b2:subsec:150302}

唯一性也可以直接从能量等式看出。若 \(u_1,u_2\) 具有同一数据，差 \(w\in H_0^1\) 满足 \(\int\nabla w\cdot\nabla v=0\)。取 \(v=w\)，得到 \(\nabla w=0\)；Poincaré 不等式推出 \(w=0\)。这个结论对不连通区域也成立，因为零边界约束在每个分支上排除了常数自由度。

\ProseParagraphBreak

内部方程则没有唯一性。例如任一调和函数都可加到 Poisson 方程的一个局部解上，而不改变源项。纯 Neumann 问题即使规定通量，仍允许每个连通分支上的加法常数。因而“取差再测试”最后得出梯度为零之后，还必须指出究竟由哪项边界条件或归一化排除了常数，不能直接省略最后一步。

\Needspace{6\baselineskip}
\begin{proposition}[弱比较原理]
\label{b2:prop:15-weak-comparison}
设整数 \(d\geq1\)，\(\Omega\subset\mathbb R^d\) 为非空有界开集，空间取实数域，\(H^{-1}(\Omega)=(H_0^1(\Omega))^*\)。设 \(f_1,f_2\in H^{-1}(\Omega)\)，\(u_1,u_2\in H_0^1(\Omega)\) 满足
\[
 \int_\Omega\nabla u_i\cdot\nabla v=\langle f_i,v\rangle
 \quad(v\in H_0^1(\Omega),\ i=1,2).
\]
若 \(\langle f_1-f_2,v\rangle\le0\) 对每个几乎处处非负的 \(v\in H_0^1(\Omega)\) 成立，则 \(u_1\le u_2\) 几乎处处。
\end{proposition}
\begin{proof}
令 \(w=u_1-u_2\)。第一卷命题22.2.3保证取正部保持实空间 \(H_0^1\)，命题21.3.5给出相应弱梯度公式。因此 \(w_+=\max\{w,0\}\in H_0^1\)，且 \(\nabla w_+=\mathbf1_{\{w>0\}}\nabla w\)。以它测试差方程，得
\[\lVert\nabla w_+\rVert_2^2=\langle f_1-f_2,w_+\rangle\le0.\]
零边界 Poincaré 不等式推出 \(w_+=0\)，即所需结论。所引截断结果的关键是先光滑逼近正部函数，在零水平集上用弱梯度为零处理折点，再利用正部在 \(H^1\) 中的强连续性传递零边界逼近；没有要求解逐点可微。
\end{proof}

这里的非负性还须在数据空间中解释。对非空有界开集 \(\Omega\)，实泛函 \(f\in H^{-1}(\Omega)\) 满足 \(f\geq0\)，是指
\[
 \langle f,v\rangle\geq0
 \quad\text{对所有几乎处处非负的 }v\in H_0^1(\Omega;\mathbb R).
\]
它是对测试函数作用的正性，不是给一般分布指定逐点值。若 \(f\in L^2(\Omega;\mathbb R)\) 且 \(f\geq0\) 几乎处处，已建立的 \(L^2\hookrightarrow H^{-1}\) 将配对写成可积的 \(\int fv\)，因而符合这个定义。不能只写 \(f\in L^1\) 就断言它属于 \(H^{-1}\)，也不能默认 \(fv\) 对全部 \(H_0^1\) 测试函数可积。将弱比较原理用于零源项与非负源项，得到非负数据的零边界弱解非负，与\chapref{b2:ch:13}经典最大值原理的符号一致。比较结论是几乎处处关系；尚未得到连续代表元时，不应改说“每个点都非负”。

\subsection{能量估计与 Galerkin 近似}
\label{b2:subsec:150303}

取解自身作测试，Cauchy--Schwarz 的对偶形式给出
\[\lVert\nabla u\rVert_2^2=\langle f,u\rangle
 \le\lVert f\rVert_{H^{-1}}\lVert\nabla u\rVert_2.\]
若梯度非零可约去一因子；若为零估计显然成立。因此解映射 \(H^{-1}\to H_0^1\) 连续。对于两组数据，同一计算给出
\[\lVert\nabla(u_1-u_2)\rVert_2\le\lVert f_1-f_2\rVert_{H^{-1}}.\]
特别地，平滑数据的近似会在能量范数中产生相应解的收敛，这为从光滑近似过渡到粗糙数据提供了稳定性依据。

\ProseParagraphBreak

若 \(f\in L^2\)，进一步有 \(\lVert\nabla u\rVert_2\le C_P\lVert f\rVert_2\)、\(\lVert u\rVert_2\le C_P^2\lVert f\rVert_2\)。注意区域放大时 \(C_P\) 也会变化：在 \(r\Omega\) 上换元表明 \(C_P\) 按长度尺度增长。这与二阶方程解相对于源项多出两个长度因子的事实一致，估计中的几何依赖不能无故删除。

\Needspace{6\baselineskip}
\begin{proposition}[有限维近似的最佳误差]
\label{b2:prop:15-galerkin-poisson}
设整数 \(d\geq1\)，\(\Omega\subset\mathbb R^d\) 为非空有界开集，空间取实数域，\(H^{-1}(\Omega)=(H_0^1(\Omega))^*\)。给定 \(f\in H^{-1}(\Omega)\) 与有限维线性子空间 \(V\subset H_0^1(\Omega)\)，设 \(u\in H_0^1(\Omega)\) 满足 \(\int_\Omega\nabla u\cdot\nabla v=\langle f,v\rangle\) 对全部 \(v\in H_0^1(\Omega)\) 成立。则存在唯一 \(u_V\in V\) 满足 \(\int_\Omega\nabla u_V\cdot\nabla v=\langle f,v\rangle\) 对所有 \(v\in V\) 成立，而且
\[\lVert\nabla(u-u_V)\rVert_2=\inf_{v\in V}\lVert\nabla(u-v)\rVert_2.\]
\end{proposition}
\begin{proof}
限制到 \(V\) 后型仍有界、强制，故有限维方程唯一可解。两式相减得 \(\int\nabla(u-u_V)\cdot\nabla v=0\) 对所有 \(v\in V\) 成立。对任意 \(v\in V\)，分解 \(u-v=(u-u_V)+(u_V-v)\)，两项在梯度内积下正交，故误差平方等于两项范数平方之和。最小值恰在 \(v=u_V\) 达到。
\end{proof}

近似解的误差受该有限维空间中的最佳逼近误差控制；因此，选取怎样的近似空间至关重要。若一列子空间递增且其并在 \(H_0^1\) 中稠密，最佳误差便趋零。递增性可改成对每个固定函数都能沿整个序列逼近的条件，但只有并集稠密不够，例如在良好空间之间反复插入零空间便会破坏收敛。存在性、离散方程的可解性与近似空间的逼近能力仍是三个可以分别核验的事实。

\ProseParagraphBreak

Poisson 型中，误差与近似空间真正正交，因而得到常数为一的最佳逼近。一般非对称型仍有从两个弱方程相减得到的 Galerkin 正交性，但它不能随意交换两个变量。下面的\term[cea-yinli]{Céa 引理}[Céa's lemma]用有界性和强制性代替勾股恒等式，给出准最佳逼近。

\Needspace{6\baselineskip}
\begin{proposition}[Céa 引理]\label{b2:prop:15-cea}
设 \(H\) 为实 Hilbert 空间，\(a:H\times H\to\mathbb R\) 为双线性型，存在 \(0\leq M<\infty\)、\(\alpha>0\)，使
\[
 |a(x,y)|\leq M\|x\|_H\cdot\|y\|_H,\qquad
 a(x,x)\geq\alpha\|x\|_H^2\quad(x,y\in H).
\]
给定 \(\ell\in H^*\) 及闭线性子空间 \(V\subset H\)，令 \(u\in H\)、\(u_V\in V\) 分别满足
\[
 a(u,y)=\ell(y)\quad(y\in H),\qquad
 a(u_V,v)=\ell(v)\quad(v\in V).
\]
则两解都存在且唯一，并且
\[
 \|u-u_V\|_H\leq\frac M\alpha\inf_{v\in V}\|u-v\|_H.
\]
特别地，每个有限维线性子空间均满足这里的闭性要求。
\end{proposition}
\begin{proof}
闭子空间 \(V\) 以继承的内积仍为 Hilbert 空间，限制后的型保留有界性与同一个强制常数，\(\ell|_V\) 连续。分别在 \(H\) 与 \(V\) 上使用\cref{b2:thm:15-lax-milgram}，得到两解。令 \(e=u-u_V\)，相减得到 \(a(e,v)=0\) 对所有 \(v\in V\) 成立。对任意 \(v\in V\)，因 \(v-u_V\in V\)，有
\[
 \alpha\|e\|_H^2\leq a(e,e)
 =a(e,u-v)+a(e,v-u_V)=a(e,u-v)
 \leq M\|e\|_H\cdot\|u-v\|_H.
\]
若 \(e=0\)，结论成立；否则约去 \(\|e\|_H\)，再对 \(v\in V\) 取下确界。误差关系 \(a(e,v)=0\) 对第二变量中的 \(v\in V\) 成立，这就足以得到同样适用于非对称型的估计。
\end{proof}

因此，强制型的离散解能否沿一列空间收敛，仍由这些空间能否逼近真解决定；有界性与强制性只增加了准最佳常数 \(M/\alpha\)。下面的分段仿射计算采用 Poisson 对称型，可以继续使用前一命题中常数为一的精确最佳逼近。

\ProseParagraphBreak

在区间 \((0,1)\) 上，这个结论可以算得很具体。取整数 \(N\geq2\)、均匀节点 \(x_j=jh\)（\(0\leq j\leq N\)）、\(h=1/N\)，令 \(V_h\) 为节点间分段仿射、两端为零的连续函数空间。基函数 \(\phi_j\)（\(1\leq j\leq N-1\)）在 \(x_j\) 取值一，在其余节点取值零。每个基函数只有相邻两段上的斜率非零，直接积分得
\[\int_0^1\phi_j'^2=\frac2h,\qquad
 \int_0^1\phi_j'\phi_{j+1}'=-\frac1h,\qquad
 \int_0^1\phi_j'\phi_k'=0\quad(|j-k|>1).\]
因此弱解的有限维近似由三对角线性方程组确定。矩阵的正定性并不需要另作行列式计算：任意系数向量对应的二次型就是分段仿射函数的梯度平方积分，零值只可能对应零函数。

\ProseParagraphBreak

例如取 \(f=1\)，精确弱解为 \(u(x)=x(1-x)/2\)。它在各小区间的导数平均等于两端函数值之差除以 \(h\)，而任意 \(v_h\in V_h\) 的导数在该段为常数。故 \(u\) 的节点线性插值 \(I_hu\) 满足 \(\int(u'- (I_hu)')v_h'=0\)，由唯一性即知 \(u_h=I_hu\)。在长度 \(h\) 的一段上，导数误差是以中点为零的线性函数，其平方积分为 \(h^3/12\)；函数误差则为 \((x-x_j)(x_{j+1}-x)/2\)。逐段积分并相加得到
\[\lVert u'-u_h'\rVert_2=\frac h{\sqrt{12}},\qquad
 \lVert u-u_h\rVert_2=\frac{h^2}{\sqrt{120}}.\]
这里第二个收敛速度比第一个快，来自二次函数的具体结构，不能仅从抽象存在定理读出。一般数据下的误差阶还需要解的正则性和近似空间的逼近估计；下一章提供其中的正则性工具。

\ProseParagraphBreak

弱解的误差也可以用残差直接测量。设整数 \(d\geq1\)、\(\Omega\subset\mathbb R^d\) 为非空有界开集，实值 \(f\in H^{-1}(\Omega)\)，\(u\in H_0^1(\Omega)\) 满足 \(\int_\Omega\nabla u\cdot\nabla v=\langle f,v\rangle\) 对所有实值 \(v\in H_0^1(\Omega)\) 成立。对任意实值 \(w\in H_0^1(\Omega)\) 定义 \(r=f+\Delta w\in H^{-1}\)，则 \(\langle r,v\rangle=\int\nabla(u-w)\cdot\nabla v\)。以 \(H_0^1\) 的梯度范数取对偶范数，对单位梯度范数的测试函数取上确界，Cauchy--Schwarz 给出上界，而在误差非零时取 \(v=(u-w)/\lVert\nabla(u-w)\rVert_2\) 达到该界。因此 \(\lVert r\rVert_{H^{-1}}=\lVert\nabla(u-w)\rVert_2\)。对分段仿射近似，残差可能含节点通量跳跃，不能只计算每段内为零的二阶导数；负范数恰好保留这些分布贡献。

\subsection{非齐次 Dirichlet 数据与 Neumann 问题}
\label{b2:subsec:150304}

取整数 \(d\geq1\)、非空有界开集 \(\Omega\subset\mathbb R^d\)、实值 \(f\in H^{-1}(\Omega)\) 与 \(g\in H^1(\Omega;\mathbb R)\)。令 \(u=g+w\)，其中 \(w\in H_0^1(\Omega;\mathbb R)\)。方程等价于
\[\int_\Omega\nabla w\cdot\nabla v
  =\langle f,v\rangle-\int_\Omega\nabla g\cdot\nabla v\qquad(v\in H_0^1).\]
右端属于 \(H^{-1}\)，范数不超过 \(\lVert f\rVert_{H^{-1}}+\lVert\nabla g\rVert_2\)。因此齐次理论给出唯一 \(w\)，从而给出唯一满足 \(u-g\in H_0^1\) 的解；并有
\[\lVert\nabla w\rVert_2\le\lVert f\rVert_{H^{-1}}+\lVert\nabla g\rVert_2.\]
完整 \(H^1\) 估计还要加上 \(g\) 的 \(L^2\) 范数。

\ProseParagraphBreak

若另一个延拓 \(\widetilde g\) 与 \(g\) 相差 \(H_0^1\) 函数，两者给出的仿射可行集完全相同；由唯一性得到同一个 \(u\)。所以解依赖边界类，而不依赖人为选择的延拓。若要比较两组边界数据，先选择它们差的一个延拓，再把上面的稳定性估计应用于差方程；优化延拓的范数便自然引出商空间形式的迹范数。

\ProseParagraphBreak

例如在 \(\mathbb R^d\) 的单位球上给定仿射边界数据 \(g(x)=a\cdot x+b\)，其中 \(a\in\mathbb R^d\)、\(b\in\mathbb R\)，且 \(f=0\)。由于 \(\int a\cdot\nabla v=0\) 对 \(v\in H_0^1\) 成立，齐次修正为零，故 \(u=g\)。对一般边界数据，弱方法也能构造能量解；在球上还可借\chapref{b2:ch:13}的表示公式识别这个解。不同构造的解一旦处于同一能量边界类，就由唯一性自动一致。


\ProseParagraphBreak

\noindent\textbf{Neumann 问题。}区域保持固定而边界值自由时，应换一个 Hilbert 空间，不能沿用零 Dirichlet 条件排除常数的唯一性论证。

\Needspace{6\baselineskip}
\begin{proposition}[Neumann 数据与归一化]
\label{b2:prop:15-neumann-normalized}
设整数 \(d\geq1\)，\(\Omega\subset\mathbb R^d\) 为非空有界连通 Lipschitz 区域，\(f\in L^2(\Omega;\mathbb R)\)、\(g_N\in L^2(\partial\Omega;\mathbb R)\)，边界积分取 \((d-1)\) 维面积测度（\(d=1\) 时为端点计数测度），\(\operatorname{Tr}\) 为 Sobolev 迹，求解类为 \(u\in H^1(\Omega;\mathbb R)\)。方程
\[\int_\Omega\nabla u\cdot\nabla v
 =\int_\Omega fv+\int_{\partial\Omega}g_N\operatorname{Tr}v
 \quad(v\in H^1(\Omega))\]
可解当且仅当 \(\int_\Omega f+\int_{\partial\Omega}g_N=0\)。指定 \(\int_\Omega u=0\) 后解唯一，且
\[\|u\|_{H^1}\le C_\Omega(\|f\|_2+\|g_N\|_{L^2(\partial\Omega)}).\]
\end{proposition}
\begin{proof}
测试 \(v=1\) 给出必要条件。令 \(H=\{v\in H^1(\Omega;\mathbb R):\int_\Omega v=0\}\)。积分泛函因 \(\Omega\) 测度有限而连续，故 \(H\) 是闭子空间；第一卷定理24.1.3的 Poincaré--Wirtinger 不等式适用于这里的有界连通 Lipschitz 区域，使 \(\|\nabla v\|_2\) 成为 \(H\) 上等价的完备范数。右端有界性来自
\[
 |\ell(v)|\le\|f\|_2\|v\|_2+
 \|g_N\|_{L^2(\partial\Omega)}\|\operatorname{Tr}v\|_{L^2(\partial\Omega)}
 \le C_\Omega(\|f\|_2+\|g_N\|_2)\|\nabla v\|_2,
\]
这里第一卷定理23.1.3给出有界 Lipschitz 区域上的 \(H^1\to L^2(\partial\Omega)\) 迹估计，刚才的零均值不等式再将 \(H^1\) 范数换成梯度范数。Lax--Milgram 在 \(H\) 上给出唯一解及估计。对任意 \(v\in H^1(\Omega)\)，分解 \(v=v_0+v_\Omega\)，其中 \(v_\Omega=|\Omega|^{-1}\int_\Omega v\)、\(v_0\in H\)。左端常数项为零，右端常数项因相容条件为零，所以受限测试解确实解全部测试方程。两解之差测试自身，得到梯度为零；连通性说明差为一个常数，归一化排除它。
\end{proof}

在 \((0,1)\) 上取 \(f=1\)、左端外向通量 \(g_N(0)=-1\)、右端通量 \(g_N(1)=0\)。总数据为 \(1-1+0=0\)。积分 \(-u''=1\) 得 \(u'=-x+C\)；左端条件 \(-u'(0)=-1\) 给 \(C=1\)，右端也满足。减去均值后 \(u=x-x^2/2-1/3\)。这个例子同时检验了相容条件中的加号和左端外法向的负号。若区域有多个连通分支，必须在每个分支分别相容、分别归一化，单个全局零均值无法排除所有常数自由度。

% !TeX root = ../../Book2.tex
\section{Riesz–Thorin 插值}
\label{b2:sec:1704}

线性算子可沿复参数改变输入和输出的指数，从而得到具有精确几何平均常数的插值估计。为使用这项方法，需要复变函数的最大模原理；本节把所需的条带估计单独证明。

\subsection{解析族}
\label{b2:subsec:170401}

下面采用\term[hadamard-sanxian-yinli]{Hadamard 三线引理}[Hadamard three-lines lemma]的下述版本。

\begin{lemma}[Hadamard 三线估计]\label{b2:thm:17-three-lines}
设 \(F\) 在 \(0<\operatorname{Re}z<1\) 全纯，在闭条带连续且有界，边界满足 \( |F(\mathrm it)|\leq M_0\)、\( |F(1+it)|\leq M_1\)。则 \(0<\theta<1\) 时
\[
 |F(\theta)|\leq M_0^{1-\theta}M_1^\theta.
\]
\end{lemma}
\begin{proof}
先设 \(M_0,M_1>0\)，令 \(G(z)=F(z)M_0^{z-1}M_1^{-z}\)，幂以实对数定义。对 \(\varepsilon>0\)，函数 \(G(z)e^{\varepsilon(z^2-1)}\) 在两条竖边的模不超过一；在上下水平边 \(\operatorname{Im}z=\pm R\) 上，指数的模不超过 \(e^{-\varepsilon R^2}\)，故当 \(R\) 足够大时也不超过一。对矩形应用最大模原理，再令 \(R\to\infty\)、\(\varepsilon\downarrow0\)。某个 \(M_i=0\) 时，先以 \(M_i+\delta\) 替代，令 \(\delta\downarrow0\)。
\end{proof}
一个\term[jiexi-zu]{解析族}[analytic family]是复参数下的算子族，其对测试输入、测试输出的矩阵元为全纯函数。使用三线估计时必须检查边界范数和条带内增长；“公式含有 \(z\)”不是解析族的充分条件。下面先用固定算子和变化的测试函数，避免讨论一般算子的定义域变化。

\subsection{Riesz–Thorin 定理}
\label{b2:subsec:170402}

\begin{theorem}[Riesz--Thorin]\label{b2:thm:17-riesz-thorin}
设 \(T\) 是定义在两个端点 \(L^{p_0}(X)+L^{p_1}(X)\) 上的兼容线性算子，底空间为 \(\sigma\)-有限测度空间。设 \(1\leq p_i,q_i\leq\infty\)，且 \(\|Tf\|_{q_i}\leq A_i\|f\|_{p_i}\)。若
\[
 \dfrac1p=\dfrac{1-\theta}{p_0}+\dfrac\theta{p_1},\qquad
 \dfrac1q=\dfrac{1-\theta}{q_0}+\dfrac\theta{q_1},
\]
则 \(\|Tf\|_q\leq A_0^{1-\theta}A_1^\theta\|f\|_p\)。
\end{theorem}
\begin{proof}
先取有限测度支集上的简单函数 \(f,g\)，分别归一化为 \(\|f\|_p=\|g\|_{q'}=1\)，并先讨论 \(p,q'<\infty\)。在非零值处写
\[
 f_z=\dfrac f{|f|}\,|f|^{p((1-z)/p_0+z/p_1)},\quad
 g_z=\dfrac g{|g|}\,|g|^{q'((1-z)/q_0'+z/q_1')}.
\]
零值处定义为零。它们是有限个固定简单函数的解析系数组合。于是 \(F(z)=\int(Tf_z)g_z\) 全纯；在竖边 \(i=0,1\) 上有 \(\|f_{i+it}\|_{p_i}=\|g_{i+it}\|_{q_i'}=1\)，无穷指数时对应模恒为一。端点 Hölder 得到边界常数 \(A_i\)。条带内各系数的模只依赖实部，因此 \(F\) 有界。三线估计在 \(z=\theta\) 给出双线性配对的上界。

对所有简单 \(g\) 取上确界，\(L^q\) 对偶范数公式得到 \(\|Tf\|_q\) 的估计；该公式也适用于 \(q=1,\infty\)，只需用简单函数近似符号或在接近本性上确界的有限测度集合上测试。若 \(p=\infty\)，两个源端点都是无穷，直接对固定 \(f_z=f\) 作同样论证，或者用有限端点目标指数的对偶插值。一般有限 \(p\) 输入由简单函数稠密性延拓；原端点定义与延拓在公共部分相同，因而没有另造一个不兼容的算子。
\end{proof}
对解析算子族 \(T_z\)，在相同构造中把 \(T\) 换为 \(T_z\)，只要所有测试矩阵元在条带中满足同样的增长条件，结论相同。若增长至多 \(e^{c|\operatorname{Im}z|}\)，先乘 \(e^{\delta(z-\theta)^2}\) 抑制增长，得到带有限常数的估计。这一版本将在曲面限制的解析插值证明中使用。

\subsection{Hausdorff–Young 再论}
\label{b2:subsec:170403}

\begin{corollary}[Hausdorff--Young]\label{b2:thm:17-hausdorff-young}
在本卷 \(e^{-2\pi\mathrm ix\cdot\xi}\) 的约定下，\(1\leq p\leq2\) 时
\[
 \|\widehat f\|_{p'}\leq\|f\|_p.
\]
\end{corollary}
\begin{proof}
Fourier 变换从 \(L^1\) 到 \(L^\infty\) 的范数至多一；\chapref{b2:ch:04}的 Plancherel 等式给出 \(L^2\) 等距性。以这两个兼容定义应用 Riesz--Thorin。对 \(1<p\leq2\)，变换先在 \(L^1\cap L^2\cap L^p\) 上定义，再依估计延拓；不宣称任意 \(L^p\) 函数的绝对积分逐点存在。
\end{proof}
插值给出常数一，却没有证明这在中间指数最优。Gaussian 的可计算范数已经提示最优常数问题比有界性更细。逆变换具有相同的两个端点，因而同样得到 \(L^p\) 到 \(L^{p'}\) 的估计。

\subsection{插值思想}
\label{b2:subsec:170404}

两种插值处理的是不同信息。Marcinkiewicz 使用大小分解和分布函数，适用于极大函数这样的次线性算子，代价是通常不精确的常数。Riesz--Thorin 使用线性叠加及复参数，只需强型端点，常数按对数凸性组合。对奇异积分，常用弱 \((1,1)\) 和强 \((2,2)\) 先得到 \(1<p<2\)，再以伴随的相同估计补齐 \(p>2\)。

\ProseParagraphBreak

若一个估计含尺度参数 \(R\)，端点常数写成 \(A_i(R)\)，插值后的幂为 \(A_0(R)^{1-\theta}A_1(R)^\theta\)。例如两个端点分别增长为 \(R^a\) 和 \(R^b\)，中间指数的增长为 \(R^{(1-\theta)a+\theta b}\)。这项尺度核算必须在对 \(R\) 求和以前完成；一个在每个尺度上成立但常数不一致的估计，不能直接推出全算子有界。

% 扩写定位：b2:subsec:170404
为了看清“把输入分配到两个端点”与函数空间之间的联系，可以引入一个直接可计算的量。对一对嵌入同一线性空间的 Banach 空间 \(X_0,X_1\)，以及 \(f\in X_0+X_1\)，定义
\[
 K(t,f;X_0,X_1)\coloneqq
 \inf_{f=f_0+f_1}\bigl(\|f_0\|_{X_0}+t\|f_1\|_{X_1}\bigr),\qquad t>0.
\]
这个下确界允许分解随 \(t\) 改变。仅取固定分解通常丢失端点附近的信息。在 Lebesgue 空间的一对 \(L^1,L^\infty\) 中，它具有完整的重排描述。
\begin{proposition}\label{b2:prop:17-k-functional}
对 \(f\in L^1+L^\infty\)，有
\[
 K(t,f;L^1,L^\infty)=\int_0^t f^*(s)\,\mathrm ds.
\]
若线性算子 \(T:X_0+X_1\to Y_0+Y_1\) 在两端点的范数分别至多 \(A_0,A_1>0\)，则
\[
 K(t,Tf;Y_0,Y_1)\leq A_0K(tA_1/A_0,f;X_0,X_1).
\]
\end{proposition}
\begin{proof}
设 \(f=f_0+f_1\)。任意 \(|E|\leq t\) 都满足
\(\int_E|f|\leq\|f_0\|_1+t\|f_1\|_\infty\)。
在无原子的 Lebesgue 空间，取超水平集合，并在必要时取水平集合的一部分，便有
\(\sup_{|E|\leq t}\int_E|f|=\int_0^t f^*\)。当超水平集合或水平值不能恰好取到最优时，用逼近集合取极限即可。因此右侧不超过每个分解的代价。

反向令 \(a=f^*(t)\)，以保持相位的截断
\[
 f_1(x)=
 \begin{cases}f(x),&|f(x)|\leq a,\\a f(x)/|f(x)|,&|f(x)|>a,\end{cases}
 \qquad f_0=f-f_1.
\]
这里 \(f_1\) 有界，且层蛋糕给出
\(\|f_0\|_1=\int_0^\infty(f^*(s)-a)_+\,\mathrm ds\)。
由重排的单调性，这等于 \(\int_0^t f^*(s)\,\mathrm ds-ta\)，水平集合上的等号不影响积分。端点 \(a=0\) 直接同样处理；若遇到下确界的极限值，先取略大的截断高度再取极限。故该分解的代价就是所求积分。

对于算子估计，每个分解产生 \(Tf=Tf_0+Tf_1\)，代价至多 \(A_0\|f_0\|_{X_0}+tA_1\|f_1\|_{X_1}\)。对全部分解取下确界即可。
\end{proof}
由此可以定义 \(0<\theta<1\)、\(1\leq q\leq\infty\) 时的实插值量
\[
 \|f\|_{\theta,q;K}=
 \left(\int_0^\infty[t^{-\theta}K(t,f)]^q\frac{\mathrm dt}{t}\right)^{1/q},
\]
无穷指数取上确界。对 \(L^1,L^\infty\)，令 \(1/p=1-\theta\)，它与前面 \(L^{p,q}\) 的拟范数等价。下面验证这一说法，而不把它当作一个新的黑箱。

\ProseParagraphBreak

单调性给出 \(K(t,f)\geq t f^*(t)\)，所以 Lorentz 量不超过 \(K\) 量。反向置 \(h(s)=s^{1/p}f^*(s)\)，则
\[
 t^{-1/p'}K(t,f)=
 \int_0^1u^{-1/p}h(tu)\,\mathrm du.
\]
在乘法测度 \(\mathrm dt/t\) 中，变换 \(t\mapsto tu\) 保持范数；Minkowski 因而给出上界
\(\int_0^1u^{-1/p}\,\mathrm du\cdot\|h\|_{L^q(\mathrm dt/t)}
=p'\|f\|_{L^{p,q}}\)。
在 \(q=\infty\) 时直接取上确界也是同一证明。算子的 \(K\) 估计于是同时传递整个高度分布的控制。这里没有推出任意两个弱型端点间的最佳 Lorentz 指数；弱型端点不是已经给定的 Banach 范数端点，其使用还需另行检查。

\ProseParagraphBreak

对两个端点都给定强型的线性算子，代入上式并换变量 \(tA_1/A_0\)，得到插值常数 \(A_0^{1-\theta}A_1^\theta\)。这个乘法结构与 Riesz--Thorin 相同，而证明只使用分解、重排和 Minkowski。前者自然保持同一个中间 Lebesgue 指数，后者自然看到 Lorentz 的第二个指数；两种方法的适用条件和产生的信息因此各有侧重。

\ProseParagraphBreak

几种插值路线在同一个算子上可以互相核对。对可积核卷积，两个同指数端点由 Young 给出，Riesz--Thorin 重现整个强型区间；若核只在临界弱空间中，集合积分与大值分解更合适；若目标是让输入带有一个细于 \(L^p\) 的第二指数，则 \(K\) 泛函直接记录高度分布。选择方法的依据是手中实际掌握的端点，而不是插值定理名字的先后次序。

\ProseParagraphBreak

实际证明还应先检查共同定义域。比如 Fourier 变换同时有 \(L^1\to L^\infty\) 与 \(L^2\to L^2\) 的界，这两种对象在 \(L^1\cap L^2\) 上一致，才可以把它看成同一个待插值算子。若两端的公式只分别定义在互不相容的空间上，即使形式上写出了两个范数界，也没有一个可供插值的共同线性算子。对奇异积分，这个相容性通常由测试函数上的主值公式与平方延拓来建立。

\ProseParagraphBreak

输出是否线性同样重要。极大算子把所有尺度的绝对平均取上确界，一般不是线性的；把输入换成复幂函数后，不能期待极大函数的矩阵元对参数解析。因此本章对它采用分布函数插值。相比之下，固定的 Fourier 变换或固定卷积是线性的，解析变化发生在测试输入与测试输出中，配对仍可验证为全纯函数。假如把“绝对值”“上确界”和“积分”在解析证明中任意交换，便可能把非解析量误放入三线引理。

\ProseParagraphBreak

常数也可以作为一个有效检查。若用两端范数 \(A_0,A_1\) 推导中间界，缩放算子 \(T\mapsto cT\) 后，中间界必须线性乘以 \(|c|\)。Riesz--Thorin 的 \(A_0^{1-\theta}A_1^\theta\) 正好具有这种齐次性；Marcinkiewicz 的未优化分割可能给出较粗的常数，但仍应在适当重选高度阈值后保持齐次。若某个推导出现两个端点常数相乘、总指数大于一，通常意味着阈值没有随算子大小调整，或只是用了一个合法但很粗的非齐次上界。

\ProseParagraphBreak

插值不能恢复被端点假设遗忘的信息。两个范数界可能来自完全不同的机制，插值只利用其共同的线性或次线性结构；它不会证明核的抵消、主值存在性、最优常数或极值函数分类。例如 Hausdorff--Young 的常数一来自基本插值，而 Gaussian 测试揭示中间指数还有更细的最佳常数问题。正向上界与某类测试函数给出的下界应分别记录，只有两者相等才完成最优常数的计算。

\ProseParagraphBreak

最后，开放指数区间中的有界性并不总能通过取极限传到端点。若估计常数在 \(p\to1\) 时发散，对每个固定输入得到的范数序列就没有统一控制；即使常数不发散，也要说明输入空间与输出空间随指数变化时采用了什么收敛。极大函数的具体反例已经给出最直接的检验：它在每个 \(p>1\) 上有界，然而一个单位区间的像有不可积尾部。弱端点的提出正是为了保留这种有限而有意义的控制。

% 本章小结（不计入编号）
\begin{chaptersummary}
弱形式减少了未知函数所需的导数阶数，却没有丢掉方程：紧支集测试辨认内部散度，测试空间和边界类决定所解的边值问题，界面项记录通量平衡。负阶空间把函数源项与散度型源项放在同一连续对偶中。Lax--Milgram 用有界性构造算子，以强制性控制逆，再排除值域的正交补，得到无需对称性的存在与稳定性。

\ProseParagraphBreak

一般椭圆主部只需有界可测系数即可定义能量弱解；低阶项可能使原型失去强制性，但加上足够大的零阶移位可恢复它。去掉移位时，紧嵌入将问题化为恒等算子的紧扰动，可解性由伴随核的有限个条件决定。整个论证首先产生 \(H^1\) 解；要得到 \(H^2\)、连续梯度乃至经典解，下一章还须逐项增加系数、数据和边界的正则性。
\end{chaptersummary}


\BookExercises

各题均为本书自拟，函数空间均取实数域。题中的 \(H^{-1}(\Omega)\) 是以梯度范数赋予 \(H_0^1(\Omega;\mathbb R)\) 后的连续对偶；一维端点与点值取第一卷命题21.4.1的绝对连续代表。前几题核验空间和数据，随后比较近似、系数扰动与共振；最后一题使用选读小节的相容条件。

% 自拟：一维点源的精确对偶范数。
\Needspace{6\baselineskip}
\begin{exercise}\label{b2:ex:15-point-source}
在 \(I=(0,1)\) 上，令 \(H^{-1}(I)=(H_0^1(I;\mathbb R))^*\)，对偶范数对应 \(\|v'\|_2\)。对 \(0<s<1\)，以绝对连续代表定义 \(\langle\delta_s,v\rangle=v(s)\)。求 \(\delta_s\) 的精确对偶范数，并说明哪个单位范数函数达到对偶上确界。
\end{exercise}
\begin{solution}
令 \(G_s(x)=(1-s)x\)（\(x\le s\)）、\(G_s(x)=s(1-x)\)（\(x\ge s\)）。对 \(v\in H_0^1(0,1)\) 的绝对连续代表，零端点值给出
\[
 \int_0^1G_s'v'=(1-s)\int_0^sv'-s\int_s^1v'
 =(1-s)v(s)+sv(s)=v(s).
\]
因此 \(|v(s)|\le\lVert G_s'\rVert_2\lVert v'\rVert_2\)，且取 \(v=G_s/\lVert G_s'\rVert_2\) 达到上界。直接计算 \(\lVert G_s'\rVert_2^2=s(1-s)^2+(1-s)s^2=s(1-s)\)，所以范数为 \(\sqrt{s(1-s)}\)。
\end{solution}

% 自拟：点质量维数结论已移入正文；改求平均点源的精确能量。
\Needspace{6\baselineskip}
\begin{exercise}\label{b2:ex:15-high-dimensional-point}
设整数 \(d\ge2\)、\(0<\varepsilon<1\)，\(B=B(0,1)\subset\mathbb R^d\)，记 \(s_d=|S^{d-1}|\)。令 \(H^{-1}(B)=(H_0^1(B;\mathbb R))^*\)，对偶范数对应 \(\|\nabla v\|_2\)。在单位球上取单位质量源
\(f_\varepsilon=\frac{d}{s_d\varepsilon^d}\cdot\mathbf1_{B(0,\varepsilon)}\)。求 \(-\Delta u_\varepsilon=f_\varepsilon\)、零 Dirichlet 边界的径向弱解，并证明
\[
 \|f_\varepsilon\|_{H^{-1}}^2
 =\frac1{s_d}\left(\frac{\varepsilon^{2-d}}{d+2}
                  +\int_\varepsilon^1r^{1-d}\,\mathrm dr\right).
\]
分别写出二维与高维的发散速度，解释它与 \(f_\varepsilon\to\delta_0\) 的分布收敛是否矛盾。
\end{exercise}
\begin{solution}
径向方程为 \(-(r^{d-1}u')'=r^{d-1}f_\varepsilon\)。在零点不留附加点源，先从零积分，得到 \(u'=-r/(s_d\varepsilon^d)\) 于内球，\(u'=-1/(s_dr^{d-1})\) 于外环。令
\[
 u_\varepsilon(r)=
 \begin{cases}
 \displaystyle\frac1{s_d}\int_\varepsilon^1t^{1-d}\,\mathrm dt
            +\frac{\varepsilon^2-r^2}{2s_d\varepsilon^d},&r\leq\varepsilon,\\[3pt]
 \displaystyle\frac1{s_d}\int_r^1t^{1-d}\,\mathrm dt,&\varepsilon\leq r\leq1.
 \end{cases}
\]
函数及径向通量在分界面相接，零点附近为光滑二次函数，边界值为零；分区积分便验证弱方程。由能量唯一性，它就是所求弱解。等式 \(\langle f_\varepsilon,v\rangle=\int\nabla u_\varepsilon\cdot\nabla v\) 与 Cauchy--Schwarz 给对偶范数上界，取 \(v=u_\varepsilon/\|\nabla u_\varepsilon\|_2\) 达到上界。径向积分故给
\[
 \|f_\varepsilon\|_{H^{-1}}^2
 =s_d\int_0^1|u_\varepsilon'(r)|^2r^{d-1}\,\mathrm dr
 =\frac1{s_d}\left(\frac{\varepsilon^{2-d}}{d+2}
                  +\int_\varepsilon^1r^{1-d}\,\mathrm dr\right).
\]
二维结果为 \((2\pi)^{-1}(1/4+\log(1/\varepsilon))\)；当 \(d>2\) 时为
\(s_d^{-1}[\varepsilon^{2-d}/(d+2)+(\varepsilon^{2-d}-1)/(d-2)]\)，范数分别按 \(\sqrt{\log(1/\varepsilon)}\) 与 \(\varepsilon^{(2-d)/2}\) 发散。另一方面，对固定光滑测试函数，球内平均趋于其中心值，所以源项在分布意义下趋于点质量。这种弱得多的收敛不要求能量对偶范数有界，因而没有矛盾。
\end{solution}

% 自拟：分段材料的显式有效导热。
\Needspace{6\baselineskip}
\begin{exercise}\label{b2:ex:15-layered-material}
令 \(a=1\) 于 \((0,1/2)\)、\(a=4\) 于 \((1/2,1)\)。在实空间 \(H^1(0,1)\) 中求 \(-(au')'=0\)、\(u(0)=0,u(1)=1\) 的弱解（端点取绝对连续代表），计算左右斜率及能量 \(E(u)=\frac12\int_0^1a|u'|^2\)，并判断它是否属于 \(H^2(0,1)\)。
\end{exercise}
\begin{solution}
通量为常数 \(q\)，且 \(1=\int_0^1u'=q(1/2+1/8)=5q/8\)，所以 \(q=8/5\)。左斜率为 \(8/5\)，右斜率为 \(2/5\)，解是通过 \((0,0),(1/2,4/5),(1,1)\) 的分段仿射函数。能量 \(\tfrac12\int a|u'|^2=\tfrac12q\int u'=4/5\)。一阶导数在中点有非零跳跃，二阶分布导数含点质量，故不属于 \(L^2\)，解不在 \(H^2\)。
\end{solution}

% 自拟：Céa引理已在正文证明；改用二维模型检验误差常数的最优性。
\Needspace{6\baselineskip}
\begin{exercise}\label{b2:ex:15-cea}
设 \(H=\mathbb R^2\) 取欧氏内积，\(\beta\in\mathbb R\)，并令
\[
 \begin{gathered}
 B_\beta=\begin{pmatrix}1&-\beta\\\beta&1\end{pmatrix},\qquad
 a_\beta(x,y)=(B_\beta x)\cdot y,\\
 \ell_\beta(y)=(-\beta,1)\cdot y,\qquad
 V=\operatorname{span}\{(1,0)\}.
 \end{gathered}
\]
求使 \(|a_\beta(x,y)|\le M\cdot|x|\cdot|y|\)、\(a_\beta(x,x)\ge\alpha|x|^2\) 对全部 \(x,y\) 成立的最小 \(M\) 和最大 \(\alpha\)。求满足 \(a_\beta(u,y)=\ell_\beta(y)\) 对全部 \(y\in H\) 成立的解 \(u\)，以及满足 \(a_\beta(u_V,y)=\ell_\beta(y)\) 对全部 \(y\in V\) 成立的解 \(u_V\in V\)，并计算
\[
 \frac{|u-u_V|}{\inf_{v\in V}|u-v|}.
\]
据此说明\cref{b2:prop:15-cea}中的常数 \(M/\alpha\) 对一般非对称强制型不能统一减小；当 \(\beta\ne0\) 时，指出 Galerkin 解与欧氏正交投影的区别。
\end{exercise}
\begin{solution}
由 \(B_\beta^{\mathsf T}B_\beta=(1+\beta^2)I\)，有 \(|B_\beta x|=\sqrt{1+\beta^2}|x|\)，故最小有界常数为 \(M=\sqrt{1+\beta^2}\)；取单位 \(y=B_\beta x/|B_\beta x|\) 可达到算子范数。又 \(a_\beta(x,x)=|x|^2\)，故最大强制常数为 \(\alpha=1\)。

全空间方程等价于 \(B_\beta u=(-\beta,1)\)，所以 \(u=(0,1)\)。写 \(u_V=(s,0)\) 并用 \((1,0)\) 测试，得到 \(s=-\beta\)。因此
\[
 e=u-u_V=(\beta,1),\qquad |e|=\sqrt{1+\beta^2},\qquad
 \inf_{v\in V}|u-v|=1,
\]
最后的下确界在欧氏投影 \(v=(0,0)\) 达到。误差比正好等于 \(M/\alpha\)，所以一般 Céa 估计中的该常数确实可以达到。另一方面，\(B_\beta e=(0,1+\beta^2)\)，故 \(a_\beta(e,w)=0\) 对全部 \(w\in V\) 成立，但 \(e\cdot(1,0)=\beta\)。当 \(\beta\ne0\) 时，这一测试正交性不是欧氏正交性，\(u_V\) 也不是欧氏最佳逼近；非对称型的准最佳逼近不能改写成 Poisson 梯度内积下的等距投影结论。
\end{solution}

% 自拟：系数在L无穷范数下的稳定性。
\Needspace{6\baselineskip}
\begin{exercise}\label{b2:ex:15-coefficient-stability}
设整数 \(d\ge1\)，\(\Omega\subset\mathbb R^d\) 为非空有界开集，\(A_n,A\in L^\infty(\Omega;\mathbb R^{d\times d})\) 满足共同椭圆下界 \(\xi^{\mathsf T}A_n(x)\xi,\xi^{\mathsf T}A(x)\xi\ge\lambda|\xi|^2\)（\(\lambda>0\)），对几乎处处 \(x\) 和全部 \(\xi\in\mathbb R^d\) 成立，且 \(\lVert A_n-A\rVert_\infty\to0\)，这里矩阵取欧氏算子范数。令 \(f\in H^{-1}(\Omega)=(H_0^1(\Omega;\mathbb R))^*\)，其对偶范数对应梯度范数；\(u_n,u\in H_0^1(\Omega;\mathbb R)\) 分别弱解 \(-\operatorname{div}(A_n\nabla u_n)=f\) 与 \(-\operatorname{div}(A\nabla u)=f\)。证明 \(u_n\to u\) 于 \(H_0^1\)，并给出定量估计。
\end{exercise}
\begin{solution}
以 \(A_n\) 为左端，两方程相减得 \(\int A_n\nabla(u_n-u)\cdot\nabla v=\int(A-A_n)\nabla u\cdot\nabla v\)。取 \(v=u_n-u\)，利用椭圆性与 Hölder 不等式，得到
\[\lVert\nabla(u_n-u)\rVert_2\le\lambda^{-1}\lVert A_n-A\rVert_\infty\lVert\nabla u\rVert_2
 \le\lambda^{-2}\lVert A_n-A\rVert_\infty\lVert f\rVert_{H^{-1}}.\]
右端趋零，Poincaré 不等式再给出完整 \(H^1\) 收敛。
\end{solution}

% 自拟：Neumann归一化与相容性。
\Needspace{6\baselineskip}
\begin{exercise}\label{b2:ex:15-neumann}
令 \(\Omega=(0,1)\cup(2,3)\)，求解类为 \(H^1(\Omega;\mathbb R)\)，纯 Neumann 通量为零，即 \(\int_\Omega u'v'=\int_\Omega fv\) 对每个 \(v\in H^1(\Omega;\mathbb R)\) 成立。证明源项在两个分支分别等于 \(1,-1\) 时，虽然总积分为零，仍无弱解。若源项改为 \(f(x)=2x-1\) 于第一分支、\(f(x)=2(x-2)-1\) 于第二分支，写出全部解，并说明只要求全域零均值为何仍不唯一。
\end{exercise}
\begin{solution}
第一种数据用第一分支的示性函数测试；它在不连通开集上属于 \(H^1\)，梯度为零，右端积分为 \(1\)，所以矛盾。第二种数据在每个分支积分为零。以局部坐标 \(s\in(0,1)\) 积分 \(-u''=2s-1\) 得 \(u'=s-s^2\)，满足两端零通量；再积分得 \(u=s^2/2-s^3/3+c_j\)，\(j=1,2\)。局部多项式均值为 \(1/12\)，全域零均值只给 \(c_1+c_2=-1/6\)，留下一个自由参数。分别要求两个分支零均值才确定 \(c_1=c_2=-1/12\)。这与正文的连通假设相吻合。
\end{solution}

% 自拟：漂移项的反对称结构。
\Needspace{6\baselineskip}
\begin{exercise}\label{b2:ex:15-divergence-free-drift}
设整数 \(d\ge1\)，\(\Omega\subset\mathbb R^d\) 为非空有界开集，\(b\in W^{1,\infty}(\Omega;\mathbb R^d)\) 且 \(\operatorname{div}b=0\) 几乎处处，\(A\in L^\infty(\Omega;\mathbb R^{d\times d})\)，并对几乎处处 \(x\)、全部 \(\xi\in\mathbb R^d\) 有 \(\xi^{\mathsf T}A(x)\xi\ge\lambda|\xi|^2\)，其中 \(\lambda>0\)。令 \(H^{-1}(\Omega)=(H_0^1(\Omega;\mathbb R))^*\) 的范数对应梯度范数。证明零边界问题 \(-\operatorname{div}(A\nabla u)+b\cdot\nabla u=f\) 对任意 \(f\in H^{-1}(\Omega)\) 在 \(H_0^1(\Omega;\mathbb R)\) 中唯一弱可解，其能量估计不含 \(\lVert b\rVert_\infty\)。
\end{exercise}
\begin{solution}
对紧支集光滑 \(u\)，有 \(\int(b\cdot\nabla u)u=\tfrac12\int b\cdot\nabla(u^2)=0\)。双线性连续性和 \(H_0^1\) 密度把该等式推广到全部 \(H_0^1\)。所以型虽可能不对称，却满足 \(a(u,u)\ge\lambda\lVert\nabla u\rVert_2^2\)。有界性仍含 \(\lVert b\rVert_\infty\)，但强制常数不含它。Lax--Milgram 给出 \(\lVert\nabla u\rVert_2\le\lambda^{-1}\lVert f\rVert_{H^{-1}}\) 及唯一性。
\end{solution}

% 自拟：不同共振频率的相容条件与全体解。
\Needspace{6\baselineskip}
\begin{optionalexercise}\label{b2:ex:15-resonance}
给定 \(q\in\mathbb R\)，在 \(H_0^1(0,\pi;\mathbb R)\) 中求解零端点问题 \(-u''-4u=\sin x+q\sin2x\)。这里弱形式是 \(\int_0^\pi(u'v'-4uv)=\int_0^\pi(\sin x+q\sin2x)v\) 对全部 \(v\in H_0^1(0,\pi;\mathbb R)\) 成立。确定哪些 \(q\) 允许弱解，并在可解时写出全部解。
\end{optionalexercise}
\begin{solution}
型对称，齐次核是 \(\operatorname{span}\{\sin2x\}\)。与核配对得相容条件 \(q\int_0^\pi\sin^22x=0\)，即 \(q=0\)。此时特解为 \(-\tfrac13\sin x\)，故全部解为 \(-\tfrac13\sin x+c\sin2x\)。核的刻画可由一维分布方程先推出 \(u''=-4u\in L^2\)，再由常系数常微分方程及端点条件得到，亦可直接使用正弦展开。因此没有遗漏较粗糙的能量解。
\end{solution}
